\documentclass[titlepage]{ltjsarticle}
\title{仙台市の魅力}
\author{仙台 四郎}
\date{\today}
\setcounter{tocdepth}{4}
\begin{document}
\maketitle
\tableofcontents
\newpage
\part{仙台市について}
仙台市（Sendai city）とは、宮城県の県庁所在地であり、「杜の都」という異名で呼ばれることもある、由緒正しき街です。この文書では、宮城県仙台市の魅力について簡単にまとめてみました。

\section{仙台の名物料理}
なんといっても仙台の名物料理には、「牛タン」「笹かま」「ずんだ餅」など、有名なものがたくさんあります。もう少しマイナーなものだと、「せり鍋」「定義山の三角油揚げ」「麻婆焼きそば」など、挙げてみればキリがありません。

\section{仙台の名所}
仙台には、「青葉城址」「東北大学」「勾当台公園」など、さまざまな名所があります。豊かな自然に囲まれ、ゆったりとした時間の中でたくさんの人が日々の活動に勤しんでいるのが、仙台という街を色濃く特徴づけています。
\subsection{東北大学のキャンパス}
東北大学は、仙台市内に点在する複数のキャンパスによって構成されています。
\paragraph{片平キャンパス}
金属材料研究所、電気通信研究所などの研究所や、大学本部などが設置されています。
\paragraph{川内キャンパス}
人文社会科学系の学部が設置されています。また、全学教育もこのキャンパスで行われます。
\paragraph{青葉山キャンパス}
工学系、理学系の学部が設置されています。
\paragraph{星稜キャンパス}
医学部、歯学部、加齢医学研究所などが設置されています。

\section{ようこそ仙台へ}
まだまだ書きたいことはありますが、今日はこのへんにしておきます。皆さん、ぜひ仙台に遊びに来てくださいね。
\end{document}

